\section{Sparse extremality and exact numerical frontiers}\label{sec:examples49}
\subsection{A multiplicative extremality condition}
The alternative in Theorem~\ref{thm:duality49} also identifies active
constraints at an optimal error. This is not the usual additive
alternation condition for a linear approximation space.

\begin{proposition}[A sparse tight certificate]\label{prop:tight49}
Suppose $0<\delta_*:=2\mathcal E_2(f)<\min_e|f_e|$ and that the
full sign tensor factors. In the resulting single chamber there is
a vector $z\in\mathbb Z_+^r\setminus\{0\}$ supported on at most
$v+1$ inequalities such that
\[
 z^TC=0,\qquad \prod_i b_i(\delta_*)^{z_i}=1.
\]
The product test can be chosen nonconstant in the threshold.
Every primal optimizer makes each inequality with $z_i>0$ tight.
The feasible set in this fixed chamber and stratum is the intersection of all multiplicative-test intervals and the endpoint conditions $U_e>0$. The full optimizer also tests the stratum endpoints separately.
\end{proposition}
\begin{proof}
The rational certificate cone has finitely many extreme rays. All its
products must be at least one at $\delta_*$. If every nonconstant test were strictly
greater than one, continuity would preserve those finitely many
inequalities at a slightly smaller threshold in the same stratum;
constant tests would be unchanged. This contradicts minimality.
Choose a tight nonconstant extreme ray and use the support bound in
Theorem~\ref{thm:duality49}. For a primal feasible point the slacks
$\log b_i-C_i y$ are nonnegative, and their $z$-weighted sum is zero.
Thus every positively weighted slack vanishes. Each base in this
stratum is constant or increasing in $\delta$, so each product test
defines an interval. This also proves the final assertion.
\end{proof}

The following example has a zero-crossing threshold larger than one
of its target entries. It therefore uses the full support theorem,
not an unproved extension of the strictly active special case.

\subsection{An exact cubic optimum with no rational minimizer}
There are two positive seeds, paired with their negatives, two commands
in one actual epoch, and two final queries. Index the positive-seed
mean tensor by $(i,a,j)\in\{0,1\}^3$. For $0<\rho\le1$, prescribe
\begin{equation}\label{eq:cubic-target49}
 f_{iaj}=\rho\begin{cases}
 1/10,&i+a+j=0,\\
 2/5,&i+a+j=1,\\
 4/5,&i+a+j=2,\\
 2/5,&i+a+j=3.
 \end{cases}
\end{equation}
All negative-seed means are their negatives. This is a finite binary
numerical experiment, with all probabilities in $[1/10,9/10]$ when
$\rho=1$ and in a smaller interval when $\rho<1$.
The two commands have different prescribed effects. This particular
table is not asserted to be a two-dimensional orthogonal orbit.

\begin{theorem}[A cubic error frontier]\label{thm:cubic49}
Let $t$ be the unique real root of
\begin{equation}\label{eq:cubic49}
 500t^3-375t^2+540t-124=0.
\end{equation}
Then $13/50<t<261/1000$ and the actual optimal two-label error of
\eqref{eq:cubic-target49} is
\begin{equation}\label{eq:cubic-opt49}
 \mathcal E_2(f)=\rho t/2.
\end{equation}
A five-entry multiplicative certificate proves the lower bound, and
explicit algebraic stochastic rows attain it. If $\rho$ is positive
rational, no all-rational stochastic machine with this same profile
attains the optimum. Nevertheless, every tolerance strictly above it
admits a dyadic machine with the same boundary labels.
\end{theorem}
\begin{proof}
Set $\rho=1$ first. Taking absolute values of all factors cannot
increase the error against the nonnegative positive-seed table.
The five entries used below are $000,011,101,110,111$. Every decomposable tensor satisfies
\begin{equation}\label{eq:cubic-circuit49}
 g_{011}g_{101}g_{110}=g_{000}g_{111}^{\,2}.
\end{equation}
If its mean error is at most $\delta<4/5$, the three factors on the
left are at least $4/5-\delta$, while the two factors on the right
are at most $1/10+\delta$ and $2/5+\delta$. Necessarily
\begin{equation}\label{eq:cubic-lower49}
 (4/5-\delta)^3\le(1/10+\delta)(2/5+\delta)^2.
\end{equation}
The difference of the right side and the left side is the polynomial
in \eqref{eq:cubic49} divided by $250$. Its derivative is strictly
positive on the real line, since $1500t^2-750t+540$ has negative
discriminant and positive leading coefficient. Direct rational
substitution at $13/50$ and $261/1000$ gives opposite signs.
Consequently \eqref{eq:cubic-lower49} fails for $\delta<t$.
This argument does not require $g_{000}$ to have been active in the
threshold sign equations.

For the reverse bound put
\[
 A=1/10+t,\qquad B=2/5+t,\qquad C=4/5-t,\qquad
 u_0=A^{1/3},\quad u_1=B^{1/3}.
\]
Use the factor $(u_0,u_1)$ in each of the three modes. The resulting
tensor has values $A$, $D=A^{2/3}B^{1/3}$, $C$, and $B$ at Hamming
weights zero, one, two, and three. Indeed, $C^3=AB^2$ is exactly
\eqref{eq:cubic49}. The errors at weights zero, two, and three are
$t$. At weight one, $D^2=AC$ and the rational bounds on $t$ give
$2/5<D<1/2$, so $|D-2/5|<1/10<t$.
For completeness the weight-one calculation is
\[
 D^2=A^{4/3}B^{2/3}=A(AB^2)^{1/3}=AC,
 \qquad \frac4{25}<\frac9{25}\frac{539}{1000}<AC
 <\frac{361}{1000}\frac{27}{50}<\frac14.
\]
All powers use their positive real branches. All factors belong to $(0,1)$.
Initialize with signed coordinate $\pm u_i$, use the two command
rows $T(u_a)$ of \eqref{eq:binaryrow47}, and decode by
$(u_j,-u_j)$. This is an actual two-label machine, not just an
unconstrained tensor. Binary TV is half the mean error.
For general $\rho$, multiply one factor mode by $\rho$.
The scaled lower constraint is explicitly
\[
 (4\rho/5-\delta)^3\le(\rho/10+\delta)(2\rho/5+\delta)^2;
\]
division by $\rho^3$ proves sharpness.

Modulo seven the polynomial in \eqref{eq:cubic49} is
$3t^3+3t^2+t+2$, with values $2,2,5,1,1,2,1$ at the seven residues.
It has no root and is irreducible over $\mathbb F_7$, hence the
primitive cubic is irreducible over $\mathbb Q$.
For rational $\rho>0$, the optimal error is therefore irrational.
Here all-rational means that every initialization, transition and decoder coefficient is rational. Such a machine at this fixed finite horizon has rational
responses and rational maximum error against the rational target.
It cannot attain \eqref{eq:cubic-opt49}.
The dyadic assertion follows by approximating the finitely many
legal rows, or by Theorem~\ref{thm:grid46}. It asserts existence
\emph{below any prescribed larger tolerance}, not equality to every
larger real error.
\end{proof}

This is field sensitivity of a \emph{positive minimum error} in a
binary-query problem. It is distinct from the classical dependence
of exact nonnegative rank on the factorization field
\cite{Shitov18}. It also goes beyond the sign criterion: every nonzero
positive-seed target in \eqref{eq:cubic-target49} has the same sign.
For reference, the exact rational interval
\[
 \frac{260362898747}{10^{12}}<t<
 \frac{260362898748}{10^{12}}
\]
is certified by the signs of the cubic. The source package includes
a dyadic upper machine and a strictly violated multiplicative lower
certificate at these two endpoints. No numerical local optimizer is
used to assert the interval.

For this minimizing machine the all-word residual sums in
Theorem~\ref{thm:gram46} and \eqref{eq:sympow46} are, respectively,
\[
 S_2=2\{5t^2+3(D-2/5)^2\},\qquad
 S_4=2\{5t^4+3(D-2/5)^4\}
\]
when $\rho=1$. Their positivity detects nonexactness of the supplied
machine. It does not prove its optimality. The circuit
\eqref{eq:cubic-circuit49} proves the lower bound against every
machine; the displayed factors prove the upper bound. These are the
distinct roles of candidate verification and global certification.

\subsection{A rational orthogonal example at every horizon}
Let $A_0=I$, $A_1=\left(\begin{smallmatrix}0&1\\1&0\end{smallmatrix}\right)$,
let the positive seeds be $(4/5,3/5)$ and $(3/5,4/5)$, and include
their negatives. Request one of the two coordinates after a word of
length $N\ge1$, at signal $0<\rho\le1$.

\begin{proposition}[A magnitude obstruction for distinct orthogonal commands]
\label{prop:orthogonal49}
The minimum all-two-label error is $\rho/20$. More generally, if a
profile has at least two labels at every cut and two distinct cuts
with exactly two labels, its minimum error is $\rho/20$, regardless
of the other widths. If a cut has only one label, its minimum error
is $2\rho/5$.
\end{proposition}
\begin{proof}
Freeze the prefix and suffix outside the two narrow cuts to identity.
For cuts $s<t$, retain the macro-words $0^{t-s}$ and $10^{t-s-1}$ inside them and fix query $j=1$.
The endpoint computation has the scalar form of
Theorem~\ref{thm:scalar47}. With one fixed query its positive-seed
mean matrix is
$\left(\begin{smallmatrix}4\rho/5&3\rho/5\\
3\rho/5&4\rho/5\end{smallmatrix}\right)$.
A decomposable approximant satisfies $g_{00}g_{11}=g_{01}g_{10}$.
At mean tolerance $\delta<3\rho/5$ this forces
$(4\rho/5-\delta)^2\le(3\rho/5+\delta)^2$, hence
$\delta\ge\rho/10$. Larger tolerances already exceed that bound.
Store only the sign of the seed, use identity command rows, and give
both queries mean $\pm7\rho/10$. Every target mean differs by exactly
$\rho/10$; this attains TV error $\rho/20$ at the whole profile.
A one-label cut makes the response independent of the sign of the
seed. Identity words and a coordinate of magnitude $4\rho/5$ force
TV error $2\rho/5$; a fair coin attains it.
\end{proof}

The orthogonal example is a magnitude test beyond sign consistency,
not an assertion that its lower bound is invisible to every separate
matrix flattening. The genuine tradeoff between separately minimal
cuts remains Theorem~\ref{thm:frontier47}. Both examples use actual nonidentical commands and a positive horizon. The identity and swap commute; noncommutativity is supplied instead by Theorem~\ref{thm:orthogonal50}.

For any of these fixed profiles, perturbing the full target probability
table by at most $\eta$ in row TV changes its optimum by at most $\eta$.
The triangle inequality proves this in both directions, even when the
perturbed target is no longer antisymmetric. Thus their strict error
margins are robust statements about all machines, not exact algebraic
identities that disappear immediately under perturbation.
